\section{Lesson 3: Eliminating Districts Entirely Trades Proportionality for Local Accountability}
\label{sec:lesson3}

County-based representation (Paper E.2) eliminates the redistricting problem
by replacing drawn districts with fixed jurisdictional units whose boundaries
are established by state statute and not subject to manipulation by
congressional legislatures \citep{dellaLibera2026e2}.  The design lesson from
this experiment is that the cure for gerrymandering that county representation
offers comes with two distinct costs: reduced local accountability through the
fractional-voting mechanism, and reduced responsiveness to population change
over time.

\subsection{The Design Logic of County Representation}

County representation applies the Huntington--Hill apportionment method ---
the same method used to allocate House seats among states since 1941
\citep{balinski2001fair} --- to allocate House seats among the 3,143 US
counties and county-equivalents.  A county with population equal to $k$ times
the per-district quota receives $k$ seats; a county with population less than
one quota receives a fractional vote weight equal to its population share
of one quota.

The gerrymandering-resistance of this system is complete: county boundaries
are not drawn by anyone in connection with congressional elections.  The
allocation is mechanically determined by population.  There is no drawing
step that can be manipulated.

\subsection{Proportionality in Competitive States}

In competitive states --- where the county composition produces both Democratic-
leaning and Republican-leaning counties in roughly proportional numbers ---
county representation produces outcomes that are at least as proportional as
algorithmic single-member redistricting and often more so.  The reason is that
county-level vote shares, while not perfectly proportional, tend to be more
moderate than single-member district vote shares, because counties are larger
and more geographically heterogeneous than the compact single-member districts
that maximize sorting.

Paper E.2 finds that in states like Wisconsin, Michigan, Pennsylvania, and
Arizona --- the quintessential competitive states of the 2010s--2020s ---
county representation produces seat--vote deviations of 1--2\% versus 3--5\%
for algorithmic single-member redistricting \citep{dellaLibera2026e2}.  This
is a meaningful improvement in proportionality that comes without any drawing
step.

\subsection{Reduced Proportionality in Highly Urbanized States}

In highly urbanized states where large metropolitan counties dominate the
population, county representation reintroduces the sorting problem in a
different form.  Cook County (Chicago) would receive approximately 8 House
seats and would elect all 8 at-large.  Cook County votes approximately 75--25
Democratic; at-large election of 8 seats from a 75--25 county would produce
6 Democratic seats and 2 Republican seats under proportional allocation --- much
better than the 7--1 outcome possible under single-member sorting, but still
not reflective of the 47\% Republican county-wide vote that Republican voters
outside Chicago proper would prefer to see represented.

Moreover, the at-large representation mechanism itself is less locally
accountable than single-member districts.  A voter in Evanston, Illinois
has no specific representative accountable to Evanston under a Cook County
at-large system; she has 8 representatives nominally accountable to all of
Cook County's 5.2 million residents.  The constituent-representative bond
that single-member districts create --- even with their sorting pathology ---
is entirely absent.

\subsection{Responsiveness to Population Change}

The second cost of county representation is reduced responsiveness to population
change.  Congressional districts are redrawn after each decennial census,
allowing the map to track population shifts within states.  County boundaries,
by contrast, are fixed by state statute and change rarely.  The allocation of
House seats to counties would be recomputed after each census, but the
geographic boundaries would remain fixed.

This means that county representation automatically adjusts for inter-county
migration (a county that grows relative to the state average receives more
seats) but cannot adjust for intra-county concentration changes.  If the
population of a large county becomes more concentrated in its urban core over
a decade --- as has happened in most major metropolitan areas since 2000 ---
county representation provides no mechanism to adjust the representation of
intra-county geographic communities.

Paper E.2 simulates the 2010--2020 population shifts under county representation
and finds that the system tracks inter-state and inter-county changes reasonably
well but systematically under-represents the fastest-growing suburban communities
within large counties, which are precisely the communities most politically
competitive and most in need of representation adjustment
\citep{dellaLibera2026e2}.

\subsection{The Fractional-Vote Implementation Problem}

The institutional novelty of fractional voting is not merely a public-relations
challenge; it creates genuine governance complications.  A county whose
population entitles it to 0.7 House seats receives a representative with 0.7
votes in Congress.  If that representative votes on a bill with 217.4 yeas
and 216.6 nays (to fabricate an extreme case), the outcome depends on how
fractional votes are tabulated.  Current House rules do not accommodate
fractional votes, and any implementation would require substantial revision
of House rules and likely the Rules of the House adopted at the start of
each Congress.

The constitutional question is more fundamental: \emph{Wesberry v.\ Sanders}
\citep{wesberry1964} requires that congressional districts be as nearly equal
in population as practicable.  County populations vary by a factor of nearly
20,000 (from Los Angeles County at 10 million to Loving County, Texas at
approximately 64), a variation that is the opposite of ``as nearly equal as
practicable.''  The only way to reconcile county representation with equal
population is through fractional voting, which has no constitutional precedent
in the House of Representatives.
